\documentclass[a4paper]{article} 
\input{head}
\begin{document}


\fancyhead[C]{}
\hrule \medskip % Upper rule
\begin{minipage}{0.295\textwidth} 
\raggedright
\footnotesize
Submitted by: Pavithra Rajan \hfill\\   
NetID: pr2725\hfill\\
\end{minipage}
\begin{minipage}{0.4\textwidth} 
\centering 
\large 
CS-GY 6923: Project\\ 
\normalsize 
Machine Learning\\ 
\end{minipage}
\begin{minipage}{0.295\textwidth} 
\raggedleft

\end{minipage}
\medskip\hrule 
\bigskip

\section{Introduction}

Scaling laws for neural language models have emerged as one of the most practically useful empirical findings in modern deep learning. Kaplan et al.\cite{1} demonstrated that model loss follows a predictable power-law as a function of model size, dataset size, and compute budget, enabling practitioners to extrapolate the performance of large, expensive models from small, cheap runs. Subsequent work by Hoffmann et al.\cite{2} refined these findings into compute-optimal training recipes. However, nearly all such studies target natural language and far less is known about whether these relationships hold for structured, programmatic representations of visual content.
\vspace{5pt}

Scalable Vector Graphics (SVG) occupies an interesting middle ground between code and visual art. An SVG file is a sequence of XML-tagged drawing commands like paths, shapes, gradients, and transforms, that a renderer interprets to produce an image. This structure makes SVGs to be treated as a programming language. Yet the semantics are fundamentally visual. Recent work such as StarVector~\cite{3} has demonstrated that language models can be trained to generate and understand SVG, opening the door to automated icon synthesis, vector graphic editing, and scalable design tools. However, it remains unclear how model capacity translates to generation quality in this domain. Does loss improve at the same rate as in natural language? Do the same power-law exponents hold? Can learning rates be transferred across scales without re-sweeping?
\vspace{5pt}

This project investigates scaling laws for decoder-only Transformer language models trained exclusively on SVG data. I make two primary contributions. First, I train five GPT-style models spanning roughly two orders of magnitude in parameter count, from 853,120 (Tiny) to 85,347,072 (XL), on a curated corpus of vector icon SVGs and fit power-law curves of the form $L = a \cdot N^{-\alpha} + c$ to characterize how validation loss scales with model size. Second, I evaluate Maximal Update Parameterization ($\mu$P)\cite{4} as a mechanism for learning-rate transfer. By reparameterizing the model in $\mu$P coordinates, the optimal learning rate at small scale should remain optimal at large scale, potentially eliminating the need for expensive per-scale hyperparameter sweeps.
\vspace{1pt}

The pipeline covers the full machine learning life cycle comprising of data collection and cleaning, BPE tokenizer training, binary dataset preparation, model definition, learning-rate sweeps, scaling runs, quantitative analysis, and a qualitative evaluation of generated SVG samples. Through this end-to-end study, the aim is to characterize whether the favorable scaling properties observed in natural language carry over to the SVG domain, and whether $\mu$P provides a reliable shortcut for hyperparameter transfer in this setting.

\section{Data}

\subsection{Dataset Description and Preprocessing Pipeline}

The training corpus is assembled from four publicly released collections in the StarVector \cite{4} family of SVG datasets, all distributed via Hugging Face: \texttt{svg-icons-simple} (primary), \texttt{svg-emoji-simple}, \texttt{svg-stack-simple}, and \texttt{svg-fonts-simple}. The first two are loaded in full, while the latter two are streamed and truncated once an estimated 100M-token training budget is met. The icons set provides clean, monochromatic vector glyphs, while the emoji set adds color, gradients, and small decorative motifs; \texttt{svg-stack} contributes diversity from in-the-wild assets such as logos and UI components; and \texttt{svg-fonts} contributes typographic glyph paths. This blend keeps the corpus dominated by icon-style SVGs while exposing the model to a wider stylistic distribution.
\vspace{4pt}

The preprocessing pipeline is intentionally aggressive, since SVGs in the real-world contain large amounts of vendor-specific markup, embedded comments, and high-precision coordinates that inflate the vocabulary without aiding generation. Each file is passed through an \texttt{lxml} parser that strips comments, processing instructions, and DOM subtrees rooted at \texttt{<metadata>}, \texttt{<title>}, \texttt{<desc>}, and \texttt{<defs>}. Empty namespace declarations are then collapsed via \texttt{etree.cleanup\_namespaces}, all numeric literals are rounded to one decimal place (e.g.\ \texttt{12.345678} $\rightarrow$ \texttt{12.3}), and runs of whitespace are normalized to a single space. SVGs that fail to parse, or whose cleaned form falls outside the $[50,\,8192]$-character window, are discarded. The lower bound rejects degenerate stubs (empty paths, single rects) and the upper bound matches roughly the 512-token context window used in training, after the BPE compression ratio is taken into account.
\vspace{4pt}

Cleaned SVGs are written to a single \texttt{all\_svgs.jsonl} artifact, which serves as input to both tokenizer training and the binary packing stage. The corpus is then tokenized end-to-end, every sequence whose token length falls outside $[4,\,513]$ is discarded, and the remaining sequences are deterministically shuffled with seed 42 and split 98/1/1 into train, validation, and test partitions \emph{at the SVG level} so that no document straddles a split boundary. Each partition is concatenated into a flat \texttt{uint16} NumPy array (\texttt{train.npy}, \texttt{val.npy}, \texttt{test.npy}) following the nanoGPT~\cite{5} layout, and the dataloader samples random contiguous windows of length \texttt{block\_size}=512 from these arrays. As every document is wrapped with \texttt{<bos>}~\dots~\texttt{<eos>}, the model receives a clear in-band signal of document boundaries even when a sampled window straddles two SVGs.

\subsection{Tokenization Scheme}

Tokenization uses Byte-Pair Encoding (BPE) with a 4{,}096-token vocabulary, a byte-level pre-tokenizer (so any UTF-8 character can be represented and \texttt{<unk>} is essentially never emitted), and a \texttt{TemplateProcessing} post-processor that wraps every encoded sequence with the special tokens \texttt{<bos>}~\dots~\texttt{<eos>}. The vocabulary also reserves \texttt{<pad>} and \texttt{<unk>}. Their integer IDs are persisted in \texttt{token\_config.json} for use by the dataloader and generator. A 4{,}096-entry vocabulary sits between the 1{,}024-token character-level extreme and the 8{,}192-token wordpiece end of the spectrum. It is large enough to capture the most common path-command motifs, but small enough that the embedding matrix remains a small fraction of total parameters at every model scale.
\vspace{4pt}

The merges learned by the trainer reveal the structure of the SVG ``language''. The very first merges fuse the leading space character (encoded as \texttt{G} in byte-level BPE) with single digits and with the path commands \texttt{L}, \texttt{C}, and \texttt{M}, reflecting how SVG path data is dominated by sequences of the form \texttt{L\,x\,y} or \texttt{C\,x\textsubscript{1}\,y\textsubscript{1}\,x\textsubscript{2}\,y\textsubscript{2}\,x\,y}. Shortly after, full attribute keywords emerge as single tokens (\texttt{Gstroke}, \texttt{Gfill}, \texttt{Gwidth}), followed by frequent integer prefixes (\texttt{G12}, \texttt{G13}, \texttt{G14}, \dots), and finally rarer multi-digit coordinates seen in the larger viewboxes of \texttt{svg-stack} (\texttt{G2343}, \texttt{1928}). As an example, the below icon
\begin{figure}[h]
    \centering
    \begin{minipage}{0.65\linewidth}
        {\small\begin{verbatim}
<svg xmlns="http://www.w3.org/2000/svg"
     viewBox="0 0 24 24">
<path d="M12 2L2 7l10 5 10-5-10-5z"/></svg>
        \end{verbatim}}
    \end{minipage}%
    \begin{minipage}{0.3\linewidth}
        \centering
        \includegraphics[width=0.7\linewidth]{diamond.png}
    \end{minipage}
    \caption{An SVG of a rhombus shape}
    \label{fig:placeholder}
\end{figure}
\noindent encodes to roughly 30 tokens including \texttt{<bos>}/\texttt{<eos>}. This is the underlying reason of using a corpus-wide compression of $\sim$3 characters per token. As noted in the estimate, this is denser than OpenAI's English benchmark of $\sim$4 chars/token because SVG contains more numerals and structural punctuation that resist merging.

\subsection{Statistics and Rendered Examples}

After cleaning, the raw JSONL contains $774{,}688$ SVGs totalling $804.6$M characters (avg.\ $1{,}038$ chars per file). Tokenization, the $[4,\,513]$-token length filter, and the document-level split together yield the partitions in Table~\ref{tab:data_splits}; the training set comprises $145.6$M tokens, comfortably above the $\sim$100M target. Figure~\ref{fig:seq_hist} shows the token-length distribution after filtering. The corpus is bimodal between the 128--256 and 384--512 buckets, reflecting the contrast between simple monochrome icons (\texttt{svg-icons}) and denser multi-path glyphs from \texttt{svg-stack} and \texttt{svg-fonts}. Only $975$ of the $\sim$435K candidate documents exceed the 512-token cap and are dropped. Average sequence length after filtering is $342$ tokens.

\begin{table}[h]
\centering
\small
\begin{tabular}{lrrr}
\toprule
Split & SVGs & Tokens & Share \\
\midrule
Train      & 425{,}063 & 145{,}573{,}946 & 98.0\% \\
Validation & 4{,}337   & 1{,}482{,}782   & 1.0\%  \\
Test       & 4{,}337   & 1{,}474{,}488   & 1.0\%  \\
\bottomrule
\end{tabular}
\caption{Final dataset partitions after BPE tokenization, length filtering, and 98/1/1 split. Token counts include \texttt{<bos>}/\texttt{<eos>} markers.}
\label{tab:data_splits}
\end{table}

\begin{figure}[h]
\centering
\includegraphics[width=0.85\linewidth]{seq_length_hist.png}
\caption{Token-length distribution of the cleaned SVG corpus. Sequences exceeding the 512-token context window ($>$512) are filtered out; the remaining $\sim$432K SVGs concentrate in the 256--512 range.}
\label{fig:seq_hist}
\end{figure}

Figure~\ref{fig:svg_examples} shows a small grid of rendered training examples.

\begin{figure}[h]
\centering
\includegraphics[width=\linewidth]{svg_examples_grid.png}
\caption{Representative SVG training examples. Top row: monochrome stroked icons from \texttt{svg-icons-simple}. Bottom row: color glyphs and multi-path documents from \texttt{svg-emoji-simple} and \texttt{svg-stack-simple}.}
\label{fig:svg_examples}
\end{figure}

\section{Methods}
\section{Results}
\section{Discussion}
\section{Conclusion}
\section{Appendices}

\begin{thebibliography}{00}

\bibitem[1]{1} J. Kaplan, S. McCandlish, T. Henighan, T. B. Brown, B. Chess, R. Child, S. Gray, A. Radford, J. Wu, and D. Amodei, ``Scaling Laws for Neural Language Models,'' \textit{arXiv preprint arXiv:2001.08361}, 2020.

\bibitem[2]{2} J. Hoffmann, S. Borgeaud, A. Mensch, E. Buchatskaya, T. Cai, E. Rutherford, D. de Las Casas, L. A. Hendricks, J. Welbl, A. Clark, et al., ``Training Compute-Optimal Large Language Models,'' \textit{arXiv preprint arXiv:2203.15556}, 2022.

\bibitem[3]{3} G. Yang, E. J. Hu, I. Babuschkin, S. Sidor, X. Liu, D. Farhi, N. Ryder, J. Pachocki, W. Chen, and J. Gao, ``Tensor Programs V: Tuning Large Neural Networks via Zero-Shot Hyperparameter Transfer,'' \textit{arXiv preprint arXiv:2203.03466}, 2022.

\bibitem[4]{4} J. A. Rodriguez, S. Agarwal, I. H. Ben Mhenni, E. Belilovsky, E. Oyallon, and S. E. Kahou, ``StarVector: Generating Scalable Vector Graphics Code from Images,'' \textit{arXiv preprint arXiv:2312.11556}, 2023.

\bibitem[5]{5} A. Karpathy, ``nanoGPT,'' \url{https://github.com/karpathy/nanoGPT}, 2022.

\end{thebibliography}

\end{document}
